\documentclass{article}
\input{../../../../lib.sty}

\title{\texttt{NoiseThresholdPrivacyMap<L01InfDistance<AbsoluteDistance<RBig>>, Approximate<MaxDivergence>> for ZExpFamily<1>}}
\author{\MichaelShoemateAuthor}
\date{}

\begin{document}

\maketitle

\contrib
Proves soundness of the implementation of \rustdoc{measurements/noise_threshold/trait}{NoiseThresholdPrivacyMap} for \texttt{ZExpFamily<1>} in \asOfCommit{mod.rs}{f5bb719}.

\section{Hoare Triple}
\subsection*{Precondition}
\subsubsection*{Compiler-verified}
\texttt{NoiseThresholdPrivacyMap} is parameterized as follows:
\begin{itemize}
    \item \texttt{MI}, the input metric, is of type \rustdoc{metrics/type}{L01InfDistance}\texttt{<AbsoluteDistance<RBig>>}
    \item \texttt{MO}, the output measure, is of type \rustdoc{measures/struct}{Approximate}\texttt{<MaxDivergence>}
\end{itemize}

\subsubsection*{Caller-verified}
None

\subsection*{Pseudocode}
\lstinputlisting[language=Python,firstline=2,escapechar=|]{./pseudocode/NoiseThresholdPrivacyMap_for_ZExpFamily1.py}

\subsection*{Postcondition}
\begin{theorem}
    Given a distribution \texttt{self},
    returns \texttt{Err(e)} if \texttt{self} is not a valid distribution.
    Otherwise the output is \texttt{Ok(privacy\_map)}
    where \texttt{privacy\_map} observes the following:

    Define \texttt{function(x)} as a function that omits entries with $v_i = 0$,
    then updates each remaining pair $(k_i, v_i + Z_i)$,
    where $Z_i$ are iid samples from \texttt{self},
    and retains pairs for which $v_i + Z_i > T$ in the positive-threshold direction,
    or $v_i + Z_i < -T$ in the negative-threshold direction, where $T = |\texttt{threshold}|$.
    The ordering of returned pairs is independent from the input ordering.

    For every pair of elements $x, x'$ in \texttt{MapDomain<AtomDomain<TK>, AtomDomain<IBig>{}>},
    and for every pair (\texttt{d\_in}, \texttt{d\_out}),
    where \texttt{d\_in} has the associated type for \texttt{input\_metric} and \texttt{d\_out} has the associated type for \texttt{output\_measure},
    if $x, x'$ are \texttt{d\_in}-close under \texttt{input\_metric}, \texttt{privacy\_map(d\_in)} does not raise an exception,
    and $\texttt{privacy\_map(d\_in)} \leq \texttt{d\_out}$,
    then \texttt{function(x)}, \texttt{function(x')} are \texttt{d\_out}-close under \texttt{output\_measure}.
\end{theorem}

\begin{proof}
    Line \ref{line:noise-privacy-map} rejects \texttt{self} if \texttt{self} does not represent a valid distribution,
    satisfying the error conditions of the postcondition.

    We now construct the privacy map.
    Both the \texttt{l1} and \texttt{li} sensitivity can be floored,
    as neighboring integer datasets always differ in whole integer increments.
    This doesn't affect the check on line \ref{line:li-check} to ensure \texttt{li} is non-negative.

    There are two ways to bound the l1 sensitivity given the three bounds:
    \begin{enumerate}
        \item The \texttt{l1} bound directly.
        \item Define $x \sim x'$ as $||x - x'||_0 \leq \texttt{l0}$ and $||x - x'||_\infty \leq \texttt{li}$. Then
        \begin{equation}
            \max_{x \sim x'} ||x - x'||_1 = \max{x \sim x'} \sum^n_i |x_i - x'_i| \leq \texttt{l0} \cdot \texttt{li}
        \end{equation}
    \end{enumerate}
    Line \ref{line:l1} updates \texttt{l1} to the tighter of these bounds.

    \texttt{li} similarly has multiple bounds:
    \begin{enumerate}
        \item The \texttt{li} bound directly.
        \item Define $x \sim x'$ as $||x - x'||_1 \leq \texttt{l1}$. Then
        \begin{equation}
            \max_{x \sim x'} ||x - x'||_\infty = \max{x \sim x'} \max_i |x_i - x'_i| \leq \max_i [0, \cdots, \texttt{l1}, \cdots, 0] = \texttt{l1}
        \end{equation}
    \end{enumerate}
    Line \ref{line:li} updates \texttt{li} to the tighter of these bounds.

    Now the \texttt{l1} sensitivity is zero if any of the three bounds are zero,
    due to the tightening of the \texttt{l1} bound on line \ref{line:l1}.
    Line \ref{line:zero-sens} then checks for zero sensitivity to return a privacy loss of zero epsilon, zero delta,
    because all neighboring datasets have identical nonzero-valued pairs after omitting zeros, and the ordering is randomized.

    Otherwise sensitivity is nonzero, so if the scale is zero,
    the privacy loss is unbounded, shown on line \ref{line:zero-scale}.

    Now that the edge cases are handled, all sensitivities and scales are finite and strictly positive.

    The mechanism that adds noise to the values in the hashmap, as described in the conditions of the postcondition of \rustdoc{measurements/noise_threshold/trait}{NoiseThresholdPrivacyMap},
    matches the mechanism that adds noise to a vector, as described in the conditions of the postcondition of \rustdoc{measurements/noise/trait}{NoisePrivacyMap}.
    Therefore releasing the values incurs the privacy loss computed by \texttt{noise\_privacy\_map} as defined on line \ref{line:epsilon},
    The privacy map also ensures that \texttt{l1} is non-negative.

    We now focus on computing the remaining privacy loss parameter delta.
    To ensure that the distance to instability remains positive,
    the threshold must be at least as large as the sensitivity, as checked on line \ref{line:threshold-check}.

    It suffices to analyze the positive-threshold direction.
    For a negative threshold $-T$, the release condition $c_j + Z_j < -T$
    is equivalent to $-c_j - Z_j > T$.
    Since the noise distribution is symmetric, $-Z_j$ has the same distribution as $Z_j$.
    Negation also preserves all input sensitivity bounds, so the privacy bound is unchanged.

    On either side of an adjacent pair, let $a$ be the magnitudes of the values of the keys absent on the other side.
    By symmetry of the noise and monotonicity of its tail, the probability of releasing any such key is at most
    \[
       1-\prod_{i:a_i>0}F(T-a_i).
    \]
    The postcondition of \texttt{threshold\_delta} bounds this probability using the resulting $L_0,L_1,L_\infty$ bounds.
    The same bound holds in the reverse direction.
    Conditioned on releasing no unstable key, the common-key noise mechanism is unchanged by independence.
    Its privacy guarantee therefore yields the claimed approximate guarantee.
\end{proof}

\bibliographystyle{alpha}
\bibliography{ref}

\end{document}
